\documentclass[11pt, a4paper]{article}

% ==================== Packages ====================
\usepackage[utf8]{inputenc}
\usepackage[T1]{fontenc}
\usepackage{geometry}
\geometry{left=2.54cm, right=2.54cm, top=2.54cm, bottom=2.54cm} % Standard 1-inch margins
\usepackage{setspace}
\onehalfspacing % Scribbr recommends 1.5 spacing
\usepackage{amsmath, amsfonts, amssymb}
\usepackage{graphicx}
\usepackage{hyperref}
\usepackage{booktabs}
\usepackage{enumitem}
\usepackage{titlesec}
\usepackage{float}
\usepackage[style=apa, backend=biber]{biblatex} % APA 7th
\addbibresource{references.bib}

\hypersetup{colorlinks=true, linkcolor=black, citecolor=blue, urlcolor=blue}

% ==================== Title Page ====================
\title{\textbf{Building Structured Retrieval and Generation Systems for Software Repositories: A Graph-Augmented Approach}}
\author{\textbf{[Your Name]} \\ [Your Department/University] \\ Supervisor: [Supervisor's Name]}
\date{\today}

\begin{document}

\maketitle
\tableofcontents
\newpage

% ==================== 1. Introduction ====================
\section{Introduction}
Large Language Models (LLMs) have demonstrated remarkable proficiency in function-level code generation. However, they face significant hurdles when applied to real-world, enterprise-level software development. Modern software projects are not isolated snippets but complex networks of hundreds of files involving class hierarchies, cross-file interface invocations, and implicit data flows. This research focuses on bridging the gap between "flat-text" retrieval and the structural complexity of software repositories to enable more accurate and efficient repository-level code generation.

% ==================== 2. Research Aims and Objectives ====================
\section{Research Aims and Objectives}
The primary aim of this research is to develop a structured retrieval framework that preserves code semantics and dependency relationships, thereby overcoming the limitations of traditional vector-based retrieval. Specifically, this project seeks to:
\begin{itemize}
    \item Develop an incremental parsing pipeline using \textit{Tree-sitter} to construct repository-scale knowledge graphs that capture both syntactic and semantic dependencies.
    \item Design a hybrid retrieval algorithm that combines dense semantic embeddings with graph-based structural expansion to extract minimal yet logically complete cross-file context.
    \item Evaluate the framework's impact on LLM functional correctness (Pass@k) and token efficiency compared to zero-shot, BM25, and standard dense retrieval baselines.
\end{itemize}

% ==================== 3. Literature Review ====================
\section{Literature Review}


% ==================== 4. Identified Gaps and Difficulties ====================
\section{Identified Gaps}
Current Retrieval-Augmented Code Generation systems suffer from two critical technical gaps when handling repository-level tasks:

\begin{itemize}
    \item \textbf{Gap 1: Semantic Loss in "Flat-Text" Retrieval.} Existing RAG systems predominantly rely on dense vector retrieval, which treats code as unstructured text. This approach ignores the inherent structural properties of code, such as ASTs and Call Graphs. Consequently, when generation requires deep-nested dependencies (e.g., second- or third-order function calls), vector-based retrieval often fails to recall the necessary context, leading to incomplete or incorrect code generation.

    \item \textbf{Gap 2: Contextual Overload and "Lost-in-the-Middle."} To compensate for retrieval inaccuracies, current engineering practices often brute-force entire files or multiple modules into the LLM's context window. This not only incurs high token costs and inference latency but also leads to attention dilution, where the model becomes distracted by redundant code, resulting in severe hallucinations and degraded generation quality.
\end{itemize}

% ==================== 5. Research Questions ====================
\section{Research Questions}
This study aims to address the following three core research questions:
\begin{enumerate}[label=RQ\arabic*:]
    \item How can heterogeneous and massive code repositories be transformed into structured dependency graphs while maintaining low parsing latency and high error tolerance?
    \item How can an efficient graph traversal and scoring mechanism be designed to precisely extract the most logical and minimal "cross-file context subgraph" for a specific query?
    \item To what extent does graph-augmented retrieval improve functional correctness (Pass@k) and reduce contextual overhead compared to traditional vector-based retrieval?
\end{enumerate}

% ==================== 6. Research Design and Methods ====================
\section{Research Design and Methods}
This study adopts a design-science research methodology, consisting of three core phases that directly correspond to the research questions outlined above.

\subsection{Phase 1: AST Extraction and Graph Construction (Addressing RQ1)}
We will utilize the \textit{Tree-sitter} incremental parsing framework to model the repository as a directed property graph $G = (V, E)$. Vertices $V$ represent code entities (functions, classes, methods) with associated dense embeddings $v_{dense}$, and edges $E$ represent explicit logical dependencies such as \texttt{CALLS}, \texttt{INHERITS\_FROM}, and \texttt{IMPLEMENTS}. A symbol resolution engine will be implemented to ensure cross-file link accuracy and handle partial or syntactically incomplete code with high error tolerance.

\subsection{Phase 2: Hybrid Retrieval and Context Pruning (Addressing RQ2)}
We propose a hybrid retrieval algorithm combining "Semantic Anchoring" with "Structural Expansion":

\begin{enumerate}
    \item \textbf{Semantic Anchoring:} Retrieve Top-$K$ seed nodes from a vector database (e.g., FAISS) based on cosine similarity $\text{Sim}_{dense}(E(q), E(v))$ between the query embedding and node embeddings.
    \item \textbf{Structural Expansion:} Execute a bounded $N$-hop traversal from seed nodes along dependency edges. The relevance score $S(v|q)$ for each candidate node is calculated as:
    \[ S(v|q) = \alpha \cdot \text{Sim}_{dense}(E(q), E(v)) + \beta \cdot \max_{u \in V_{seed}} (\gamma^{d(v, u)}) \]
    where $\alpha$ and $\beta$ are weighting coefficients, $\gamma \in (0,1)$ is a distance decay factor, and $d(v, u)$ is the shortest path length between nodes $v$ and $u$.
    \item \textbf{Linearization:} The resulting subgraph is topologically sorted and formatted into structured Markdown/XML tags for efficient LLM consumption.
\end{enumerate}

\subsection{Phase 3: Empirical Evaluation (Addressing RQ3)}
The system will be evaluated on authoritative repository-level datasets, including \textbf{CrossCodeEval} and \textbf{SWE-bench (Lite)}. We will compare our Graph-RAG method against the following baselines:
\begin{itemize}
    \item \textbf{Zero-shot LLMs:} Direct generation without any retrieval augmentation.
    \item \textbf{BM25:} Sparse lexical retrieval as a traditional information retrieval baseline.
    \item \textbf{Standard Dense RAG:} Pure vector-based retrieval without structural expansion.
\end{itemize}
Evaluation metrics include:
\begin{itemize}
    \item \textbf{Functional Correctness:} Pass@1 and Pass@10 through execution-based testing, supplemented by CodeBLEU for surface-level similarity.
    \item \textbf{Efficiency:} Token Reduction Rate, measuring the percentage reduction in prompt tokens compared to full-repository baselines.
\end{itemize}

% ==================== 7. Significance and Contribution ====================
\section{Significance and Contribution}
This research contributes to the field of AI for Software Engineering (AI4SE) by providing a scalable solution for repository-level reasoning. By shifting from unstructured text retrieval to structural dependency alignment through graph-augmented retrieval, this project promises to (1) significantly improve the functional correctness of LLM-generated code in complex, multi-file settings, and (2) reduce the cost and latency of LLM inference by delivering precise, minimal-context prompts rather than brute-force file injection. The proposed framework has the potential to bridge the gap between laboratory benchmarks and real-world enterprise software development.
% ==================== 8. Project Timeline ====================
\section{Project Timeline}
The proposed research is designed to be completed within a four-month period. Table~\ref{tab:timeline} presents the week-by-week breakdown across the three phases, with overlapping milestones to reflect the iterative nature of system development and evaluation.

\begin{table}[H]
\centering
\caption{Four-Month Project Timeline}
\label{tab:timeline}
\begin{tabular}{@{}p{2.2cm} p{2.4cm} p{8.5cm}@{}}
\toprule
\textbf{Timeframe} & \textbf{Phase} & \textbf{Key Activities \& Deliverables} \\
\midrule
Month 1, Weeks 1--2 & Phase 1 & \textbf{Literature review \& environment setup.} Finalize dataset selection (CrossCodeEval, SWE-bench Lite); set up Tree-sitter parsing pipeline for Python; implement AST node extraction and symbol table construction. \\
\addlinespace
Month 1, Weeks 3--4 & Phase 1 & \textbf{Graph construction.} Implement directed property graph builder with \texttt{CALLS} and \texttt{INHERITS\_FROM} edge types; integrate dense embedding generation (e.g., CodeBERT) and FAISS index construction; produce preliminary graph statistics report. \\
\midrule
Month 2, Weeks 1--2 & Phase 2 & \textbf{Hybrid retrieval — Semantic Anchoring.} Implement cosine-similarity-based seed retrieval via FAISS; conduct ablation on Top-K values; document baseline retrieval recall metrics. \\
\addlinespace
Month 2, Weeks 3--4 & Phase 2 & \textbf{Hybrid retrieval — Structural Expansion.} Implement $N$-hop bounded graph traversal with decay-weighted scoring function $S(v|q)$; tune hyperparameters $\alpha$, $\beta$, $\gamma$; develop topological-sort linearizer; produce end-to-end retrieval prototype. \\
\midrule
Month 3, Weeks 1--2 & Phase 3 & \textbf{Baseline evaluation.} Run zero-shot LLM, BM25, and standard dense RAG baselines on both datasets; collect Pass@k and token-usage statistics; identify failure cases for qualitative analysis. \\
\addlinespace
Month 3, Weeks 3--4 & Phase 3 & \textbf{Graph-RAG evaluation.} Execute full Graph-RAG pipeline on evaluation datasets; compute Pass@1, Pass@10, CodeBLEU, and Token Reduction Rate; conduct statistical significance tests between conditions. \\
\midrule
Month 4, Weeks 1--2 & All phases & \textbf{Ablation studies \& refinement.} Perform ablation on hop-count limit, decay factor $\gamma$, and embedding model choice; refine scoring function based on error analysis; rerun critical experiments. \\
\addlinespace
Month 4, Weeks 3--4 & All phases & \textbf{Write-up \& dissemination.} Consolidate results, figures, and statistical analyses; draft and revise the final thesis/manuscript; prepare supplementary materials (code repository, replication package). \\
\bottomrule
\end{tabular}
\end{table}

\noindent The schedule allocates approximately 25\% of the total duration to graph infrastructure (Phase 1), 25\% to retrieval algorithm design (Phase 2), 25\% to empirical evaluation (Phase 3), and 25\% to refinement and write-up. A one-week buffer is implicitly reserved at the end of Month~4 for unforeseen delays.

% ==================== 9. Risk Assessment and Mitigation ====================
\section{Risk Assessment and Mitigation}
Table~\ref{tab:risks} summarizes the principal risks identified for this project and the corresponding mitigation strategies.

\begin{table}
\centering
\caption{Risk Assessment and Mitigation Strategies}
\label{tab:risks}
\begin{tabular}{@{}p{5cm} p{2cm} p{6.7cm}@{}}
\toprule
\textbf{Risk} & \textbf{Severity} & \textbf{Mitigation Strategy} \\
\midrule
\textbf{R1: Parsing failures on heterogeneous codebases.} Tree-sitter may fail on non-standard language dialects, preprocessor macros, or syntactically incomplete code prevalent in real-world repositories. & High & Implement a \textit{best-effort} parsing mode that gracefully degrades to regex-based fallback extraction for unparseable files. Validate on a curated subset of repositories with known parsing challenges (e.g., SWE-bench's diverse codebases) during Month~1 to surface issues early. \\
\hline
\textbf{R2: Graph explosion during multi-hop traversal.} Unbounded structural expansion from seed nodes may retrieve an excessively large subgraph, negating the token-efficiency goals of the project. & High & Enforce a strict hop-count limit $N \leq 3$ based on empirical findings from prior work \cite{liu2025enhancing}. Additionally, apply a relevance-score threshold to prune low-weight nodes post-expansion. Monitor subgraph size distributions during Month~2 prototyping to calibrate thresholds. \\
\hline
\textbf{R3: Embedding model misalignment with code semantics.} General-purpose text embedding models may produce poor semantic representations for code-specific constructs, undermining the anchoring stage. & Medium & Evaluate both general-purpose (e.g., text-embedding-3) and code-specific models (e.g., CodeBERT, UniXcoder). Select the best-performing model via a pilot study on CodeSearchNet during Month~1. \\
\hline
\textbf{R4: Dataset contamination.} LLM-based baselines may have memorized evaluation benchmarks (CrossCodeEval, SWE-bench) during pre-training, inflating baseline performance and masking the true benefit of Graph-RAG. & Medium & Cross-reference evaluation samples against known training corpora where possible. Supplement quantitative metrics with qualitative case studies on newly created, held-out repository tasks to validate generalizability. \\
\hline
\textbf{R5: LLM API cost and rate limits.} Extensive evaluation across multiple baselines and ablation configurations may incur significant API costs or hit rate limits during Month~3 and Month~4. & Medium & Use open-weight models (e.g., DeepSeek-Coder, CodeLlama) for development and debugging iterations. Reserve paid API access for final benchmark runs only. Implement request caching to avoid redundant LLM calls for identical prompts. Budget approximately 2--3 weeks of API quota headroom. \\
\hline
\textbf{R6: Insufficient time for write-up.} The four-month scope leaves limited margin; delays in Phases~1--3 could compress the writing period and reduce thesis quality. & Medium & Adopt an \textit{incremental writing} discipline: draft the literature review and methodology sections concurrently with implementation (Months~1--2). Maintain a lab notebook with daily experiment logs to accelerate the results write-up in Month~4. Schedule weekly supervisor check-ins to ensure milestone adherence. \\
\bottomrule
\end{tabular}
\end{table}

% ==================== 8. Reference List ====================
\newpage
\printbibliography[title={Reference List}]

\end{document}
